\documentclass[11pt,a4paper]{article}
\usepackage[utf8]{inputenc}
\usepackage[margin=0.85in]{geometry}
\usepackage{amsmath,amssymb,amsfonts}
\usepackage{graphicx}
\usepackage{xcolor}
\usepackage{tcolorbox}
\usepackage{circuitikz}
\usepackage{booktabs}
\usepackage{fancyhdr}
\usepackage{hyperref}

\tcbuselibrary{skins,breakable}

\definecolor{primaryblue}{RGB}{20, 50, 110}
\definecolor{accentteal}{RGB}{0, 128, 128}
\definecolor{boxbg}{RGB}{245, 248, 252}
\definecolor{alertred}{RGB}{180, 40, 40}

\newtcolorbox{questionbox}[1]{
  colback=boxbg,
  colframe=primaryblue,
  fonttitle=\bfseries\sffamily,
  title={#1},
  arc=3mm,
  boxrule=1.2pt,
  breakable
}

\newtcolorbox{answerbox}[1][]{
  colback=green!5!white,
  colframe=green!50!black,
  fonttitle=\bfseries\sffamily,
  title={Final Result},
  arc=2mm,
  boxrule=1pt,
  #1
}

\begin{document}
\section{Test Section}
\begin{questionbox}{2083 Baishakh --- Question 1}
Find the current in 10 V source using nodal analysis.
\begin{center}
\begin{circuitikz}[american, scale=1.1, transform shape]
  % Ground rail
  \draw (0,0) node[ground]{} -- (6,0);
  % Nodes
  \draw (0,0) to[R, l=$2\,\Omega$, i>_=$i$] (0,2.5) node[left]{$v_1$}
        to[V, v<=$10\text{ V}$] (3,2.5) node[above]{$v_2$}
        to[cV, v^<=$5i$] (6,2.5) node[right]{$v_3$}
        to[R, l=$3\,\Omega$] (6,0);
  \draw (3,2.5) to[R, l=$4\,\Omega$] (3,0);
  \draw (0,2.5) -- (0,3.8) to[R, l=$6\,\Omega$] (6,3.8) -- (6,2.5);
\end{circuitikz}
\end{center}
\end{questionbox}
\end{document}
